\documentclass[
  a4paper, 
  10pt, 
  parskip=half, 
  DIV=20,
  toc=bibliography,
  usegeometry,
]{scrartcl}

% \usepackage{showframe}

% \usepackage[spanish, es-lcroman]{babel}

% \usepackage{multicol}

% \usepackage{subfiles} 
% \usepackage{tocbibind}
% \usepackage[toc]{appendix}

% \usepackage[justification=centering]{caption}
% \usepackage{subcaption}

% \usepackage[shortlabels]{enumitem}
% \usepackage{epigraph}

\usepackage{fontspec}
% switch to a monospace font supporting more Unicode characters
\setmonofont{JetBrains Mono}[
    Contextuals = Alternate,
    Ligatures = TeX,
    Scale=MatchLowercase
]

% \usepackage[outputdir=../build]{minted}
% \usemintedstyle{pastie}
\usepackage{listings}
\lstset{
    basicstyle = \ttfamily,
    columns = flexible,
}

% Use KOMA-Script's nice headers instead of fancyhdr
% \usepackage[Conny]{fncychap}

% === Math facilities ===

% Theorem environments
\usepackage{amsthm}

\newtheorem{theorem}{Theorem}[section]
% \newtheorem{proposition}[theorem]{Proposition}
% \newtheorem{corollary}[theorem]{Corollary}
% \newtheorem{lemma}[theorem]{Lemma}

\newtheorem{definition}[theorem]{Definition}
\newtheorem{example}[theorem]{Example}
\newtheorem{problem}[theorem]{Problem}

% \newtheorem{remark}[theorem]{Remark}
\newtheorem{notation}[theorem]{Notation}

\usepackage{centernot}
\usepackage{amssymb}
\usepackage{mathtools}
\usepackage{tikz-cd}

\newcommand{\funfont}[1]{\text{#1}}
\newcommand{\tyfont}[1]{\texttt{#1}}
\renewcommand{\implies}{\rightarrow}
\renewcommand{\impliedby}{\leftarrow}
\renewcommand{\mapsto}{\Rightarrow}
\newcommand{\Prop}{\mathbf{Prop}}

\usepackage{scalerel}

\DeclareMathOperator*{\Forall}{\scalerel*{\forall}{\sum}}


% === End math facilities ===

\usepackage{csquotes}

\usepackage{bookmark}
\hypersetup{
    colorlinks=true,
    % linkcolor=blue,
    % filecolor=magenta,      
    % urlcolor=cyan,
    % pdftitle={Overleaf Example},
    % pdfpagemode=FullScreen,
    }

% Los apéndices se llaman apéndices
% \renewcommand\appendixpagename{Apéndices}
% \renewcommand{\appendixtocname}{Apéndices}

% mainmatter no reinicia el contador de páginas
% \makeatletter
% \renewcommand\mainmatter{
%     \@mainmattertrue\cleardoublepage\renewcommand\thepage{\arabic{page}}}
% \makeatother

% \renewcommand{\thesubsection}{\thesection.\alph{subsection}}

% \providecommand{\keywords}[1]
% {
%   \small	
%   \textbf{\textit{Palabras clave---}} #1
% }

% define math symbol for big step arrow
\newcommand{\bigstep}{\Longrightarrow}

\title{Classic results on programming language semantics in Lean 4 (Draft)}
\author{Ricardo Maurizio Paul}
% \date{2020--2024}

\begin{document}

\maketitle

\tableofcontents

\section{Introduction}

Proof assistants such as Lean and Coq have gained significant popularity in
recent years, specifically in the field of formal verification. As software
systems become increasingly more complex, and are responsible for more
critical tasks, the need for formal verification becomes more apparent. Formal
verification is a technique that uses mathematical logic to prove the
correctness of a system, usually hardware or software. It involves using formal
methods and tools to prove that the system satisfies a set of well defined
properties that model the desired behaviour. A classic approach to formal
verification is Hoare logic, a formal system that deconstructs programs into
components each with a set of pre and post conditions, these conditions and
components, called commands, are represented as \emph{Hoare triples} of the form
\(\{P\} \ C \ \{Q\}\),
where P is the precondition, C is the command, and Q is the postcondition. By
using inference rules, Hoare logic can be used to show the correctness of a
program by proving that the postcondition of the program is satisfied given the
precondition. However, reasoning about loops and recursion can be difficult in
Hoare logic, as it requires the use of invariants to prove the termination of
the program. We will see that this is a common problem when reasoning about
programs in general, as proving the termination of a given program is a
non-trivial task, and in the most general case, is undecidable, this is known as
the famous \emph{Halting Problem}.

The inference rules for Hoare logic however are dependent on the \emph{semantics} of
the programming language. The semantics of a language are the rules that define 
the meaning of the programs, we
will see that these rules can present themselves in different forms, such as
operational semantics, denotational semantics, and axiomatic semantics. In this
text we will focus on operational semantics, and denotational semantics.
Axiomatic semantics, such as Hoare logic, are derived from operational
semantics.

These formal methods ensure \emph{correctness} of the programs, and are a more
rigorous approach to program verification than testing. Testing can only show
the presence of bugs, but not their absence, and is not exhaustive. In this way
formal methods are more suitable for critical systems, such as avionics, medical
devices, and financial systems, where the cost of failure is high. Alongside the
correctness gains, formal methods can also improve the performance of programs,
verified toolchains, such as CompCert, have been shown to produce more
efficient code than their unverified counterparts. This is due to more
aggressive optimizations being possible when the correctness of the program is
guaranteed: runtime checks can be removed, aggressive flow analysis can be
performed, and dead code can be eliminated. Some of these optimizations can be
performed by modern compilers, using static analysis, but the guarantees
provided by formal methods are stronger, as we said before we can ensure
\emph{termination} in some cases.

Proof of the relevance that program correctness has gained in recent years is
the adoption of the \emph{Rust} programming language, which has a strong type system
that ensures memory safety and thread safety, and is designed to prevent common
programming errors, such as null pointer dereferences, buffer overflows, and
data races. It ensures these properties by using a \emph{borrow checker} that
enforces strict rules on the use of memory, and by using a \emph{type system} that
ensures correctness of the programs. The borrow checker is based on
\emph{substructural logic} and \emph{linear types}, concepts that we will not cover in
this text, but that encode in the language itself the rules of the type system,
ensuring that the programs are well-typed, and that the memory is used
correctly.

\section{Proof assistants}

Proof assistants, or interactive theorem provers, are software systems that aid
in the construction and verification of mathematical proofs using computers. The
inception of proof assistants can be traced back to the 1950s and 1960s, with
the \emph{Logic Theorist}\footnote{\url{https://en.wikipedia.org/wiki/Logic_Theorist}} 
and \emph{Automath} systems. These were the
first systems to use a computer to check mathematical proofs. They were
created before the advent of modern type theory and the Curry-Howard
isomorphism, and used a different logical foundation based on \emph{natural
deduction}. Although interesting these systems were limited in scope, and
could only prove simple theorems.

In the 1970s more powerful systems were developed, such as \emph{LCF}, based on
Dana Scott's \emph{Logic of Computable Functions}, this was one of the first systems
to use proofs by induction, a feature fundamental to reason about recursively
defined functions and that we will see in this text. The creation of the
programming language ML, the father of all functional languages, is also
attributed to the development of LCF, as it was designed for writing tactics
in LCF. The modern approach of \enquote{step by step} proof construction, used by all
modern assistants, was also pioneered by LCF.\@

The 1970s also saw the development of ideas such as type theory, in particular
Martin-Löf type theory, and the Curry-Howard isomorphism, two key concepts that
are the basis of all modern proof assistants. In the 1980s this ideas were
implemented in systems such as \emph{Nuprl} and \emph{Coq}. Nuprl is one of the
first systems to use a constructive type theory at its core, were proofs are
first class objects, emphasizing their computational content and allowing
\emph{code extraction} of programs from verified proofs. Nuprl also introduced
dependent types, a powerful feature that allows the type of a value to depend on
the value itself. With this system a great deal of mathematics can be encoded
and proved.

A little later came Coq, a system based on the CIC (Calculus of Inductive
Constructions), a type theory that extends the Calculus of Constructions with
inductive types. Coq is one of the most popular proof assistants, and is widely
used in academia and industry to this day. We will see that the CIC is also the
basis of Lean, a more recent proof assistant that has gained significant
popularity in recent years.

While this advancements where made with the theories of dependent types another
line of research was being developed, the \emph{HOL} family of proof assistants,
based on \emph{Higher Order Logic}, a logic that extends first-order logic with
higher order quantifiers. These systems are based on the \emph{LCF} architecture, and
use a \emph{kernel} to ensure the correctness of the proofs. The most popular system
in this family is \emph{Isabelle, a powerful proof assistant that is as
powerful as Coq and Lean}, but with a different logical foundation.

This brings us to Lean, the proof assistant that we will use in this text. Lean
is a modern proof assistant developed by Leonardo de Moura at Microsoft Research
and based on the CIC.\@ We can see Lean as an evolution of Coq, with a more modern
syntax, and a more powerful type system, later we will discuss the differences
in more detail.

Leaving HOL aside, all these systems are based on the Curry-Howard isomorphism,
a fundamental concept that relates proofs and programs. To understand it we must
first introduce the basics of type theory.

\section{Type theory}

Type theory is a branch of mathematical logic that deals with the study of
types, an abstraction similar in some ways to the more familiar sets.
It emerged in the early 20th century to solve the foundational crisis of
mathematics, a series of problems that arose from the discovery of paradoxes in
set theory, such as Russell's paradox.

Like \emph{sets} are a primitive notion in set theory, \emph{types} are a primitive notion
in type theory, alongside \emph{functions}. This two primitives ensure that functions
are only applied to the correct types, this is known as \emph{type safety}. This
concept may be familiar to programmers, as it is the basis of static type
systems in programming languages, such as Haskell, Rust, and Scala.

As an informal introduction to types we can think of them as sets with
additional structure which dictates with functions can be applied to them. For
example, we can define a type \(\mathbb{N}\) that represents the natural numbers,
and define a function \(\texttt{add}: \mathbb{N} \to \mathbb{N} \to \mathbb{N}\)
that takes two natural numbers and returns their sum. This ensures that the
function \(\texttt{add}\) can only be applied to natural numbers, and not to any
other type.

Historically the origins of type theory can be traced back to multiple 
developments in logic, by Church, Curry, Russell, and others. But modern type
theory is based on the work of Pier Martin-Löf, who developed a type theory that
combines classic type theory with the \(\lambda\)-calculus. In this way
the \enquote{nice} properties of previous type theories are given a 
\emph{computational interpretation}, via the Curry-Howard isomorphism which
we will discuss in more detail later, that allows computers to reason about
mathematical proofs.

The type theory of Martin-Löf is \emph{intuitionistic}, that is
proofs are not just a logical argument but a method of constructing a
\emph{witness} of the truth of a proposition. This is in contrast to classical
logic, where proofs are just a sequence of logical steps that show the truth of a
proposition. In intuitionistic logic, the law of the excluded middle is not
valid, and the principle of \emph{proof by contradiction} is not allowed (with
some caveats).

\subsection{Lambda calculus}

A complete treatment on \(\lambda\)-calculus can be found in~\cite{barendregt1984lambda}, 
but for completeness we will give a brief introduction here.
The \(\lambda\)-calculus is a formal system for expressing computation based on function abstraction 
and application using variable binding and substitution. It was introduced by 
Alonzo Church in the 1930s as a way to study the foundations of mathematics and 
has since become a fundamental concept in computer science. We can think of
it as a \enquote{programming language} with only one data type, the function, and
two operations, abstraction and application.

The \(\lambda\)-calculus treats functions as \emph{rules} rather than
relationships, this stresses the computational aspect of functions, and allows
us to reason about computation in a more abstract way. This is in contrast to
set theory, where functions are seen as \emph{graphs}.

The syntax of lambda calculus is simple, it consists of variables (\(x, y, z\)),
abstraction (\(\lambda x \mapsto M\)), and application (\(M \ M\)), in 
BNF:
\[
  M \Coloneqq x \mid \lambda x \mapsto M \mid M \ M
\]
Our usage of the (\(\mapsto\)) symbol, in contrast to the traditional (\(.\))
is not accidental, and reveals a profound connection between function abstraction
and logical implication, a connection that we will explore later.

\begin{example}\label{ex:lambda_expr}
We can apply the successor function to the number 3 as
\(
(\lambda x \mapsto x + 1) \ 3
\).
\end{example}

The rules to evaluate lambda calculus expressions are known as
\emph{reduction rules} and they are:
\begin{itemize}
  \item \(\alpha\)-conversion: Allows us to rename bound variables:
    \[(\lambda x \mapsto M) \ N \rightsquigarrow M[x \coloneqq N]\]
  \item Substitution: Allows us to replace all \emph{free} occurrences of the 
  variable \(x\) in the expression \(M\) with expression \(N\) and is denoted as
  \(M[x \coloneqq N]\).
  \item \(\beta\)-reduction: Allows us to apply a function to an argument, and
    is defined in terms of substitution as:
    \[(\lambda x \mapsto M) \ N \rightsquigarrow M[x \coloneqq N]\]
  \item \(\eta\)-conversion: Expresses the concept of \emph{extensionality} of functions, 
    and states that two functions are equal if they are equal for all arguments.
    To give an example when \(x\) is not free in \(f\):
    \[(\lambda x \mapsto f \ x) \rightsquigarrow f\]
\end{itemize}

\begin{example}
  We can apply the \(\beta\)-reduction rule to the expression in
  Example~\ref{ex:lambda_expr} as:
  \[
  (\lambda x \mapsto x + 1) \ 3 \rightsquigarrow 3 + 1 
  \]
\end{example}

These reduction rules are called like that because they only go one way, from
left to right. Seeing the previous example one could think that they also
\enquote{reduce} expressions to \enquote{simpler} ones: applying the rules one
after another we would arrive to a version of the expression on which the reduction rules cannot 
be applied any more, this is called the \emph{normal form} of the expression.

\begin{definition}[Normal form]
  An expression is in normal form if it cannot be reduced further.
\end{definition}

Sadly, not all expressions have a normal form, as we will see in the next example.

\begin{definition}[\(\Omega\)-combinator]\label{def:omega_combinator}
  The \(\Omega\)-combinator is the \(\lambda\)-expression:
  \[
  \Omega \coloneqq (\lambda x \mapsto x \ x) \ (\lambda x \mapsto x \ x) 
  \]
\end{definition}

Without going into the details of the definition if the rules of the syntax and
reduction of \(\lambda\)-calculus we can't give a real \enquote{rigorous} proof
of the following theorem, but we can give an informal one.

\begin{theorem}
  The \(\Omega\)-combinator does not have a normal form.
\end{theorem}

\begin{proof}
  If we apply the \(\beta\)-reduction rule to~\ref{def:omega_combinator}
  we get the expression itself:
  \[
  (\lambda x \mapsto x \ x) \ (\lambda x \mapsto x \ x) \rightsquigarrow
  (\lambda x \mapsto x \ x) \ (\lambda x \mapsto x \ x)
  \]
\end{proof}

Since our only \enquote{data types} are functions, we need to encode numbers,
booleans, and any other data types using functions. We can use the
\emph{Church encodings}, which represents the booleans as functions that take two
arguments and return one of them. For example, we can define the boolean literals
as follows:

\begin{definition}[Church booleans]
  We can define the boolean literals as:
  \[
    \texttt{true} \coloneqq \lambda x \mapsto \lambda y \mapsto x \qquad
    \texttt{false} \coloneqq \lambda x \mapsto \lambda y \mapsto y 
  \]
\end{definition}

We can then define the logical operations \(\texttt{and}\), \(\texttt{or}\), and
\(\texttt{not}\) as functions that take two booleans and return a boolean.

\begin{definition}[Logical operators]
  We can define the logical operators as:
  \[
    \texttt{and} \coloneqq \lambda x \mapsto \lambda y \mapsto x \ y \ x \quad
    \texttt{or} \coloneqq \lambda x \mapsto \lambda y \mapsto x \ x \ y \quad
    \texttt{not} \coloneqq \lambda x \mapsto x \ \texttt{false} \ \texttt{true} 
  \]
\end{definition}

And even an if-then-else construct:

\begin{definition}[If-then-else]
  We can define the if-then-else construct as:
  \[
  \texttt{ifthen} \coloneqq \lambda b \mapsto \lambda x \mapsto \lambda y \mapsto b \ x \ y 
  \]
\end{definition}

Numbers also have a church encoding but we skip presenting it since later we 
will see a more general way to encode data types.

As we've seen the \(\lambda\)-calculus is a simple but powerful formal system
that can be used to reason about computation, it is in fact \emph{Turing complete,
meaning that it can express any computable function.} However, as we have 
shown with the \(\Omega\)-combinator, it is not very \enquote{safe} to use.

\subsection{Typed lambda calculi}

To avoid the problem introduced by non-terminating computations such as the 
\(\Omega\)-combinator we can introduce \emph{types} to the lambda calculus\footnote{
  In some type systems the \(\Omega\)-combinator can be well typed, this is not
  the case in the type systems normally used in theorem provers, such as CIC.
}. 
Types allow us to classify the expressions
of the lambda calculus and ensure that the expressions are well-typed, that is,
that the functions are applied to the correct arguments. This is known as
\emph{type safety} and is a fundamental concept in programming languages.

By adding types we also ensure that nonsensical expressions, such as applying a
boolean to an integer are not possible, and so expressions cannot get
\enquote{stuck} in a non-terminating computation.

\subsubsection{Simply typed lambda calculus}

The simplest typed lambda calculus is the \emph{simply typed lambda calculus} or
\emph{STLC}. In the STLC each variable has a type, and each function has a type
function abstraction now has the form \(\lambda x: \alpha \mapsto M\), where
\(\alpha\) is the type of the variable \(x\). The two mandatory types in the STLC
are the function type \(\alpha \to \beta\), which represents a function that
takes an argument of type \(\alpha\) and returns a value of type \(\beta\), and
the base type \(\texttt{Bool}\), which represents the booleans.

The syntax of the STLC is:
\begin{align*}
  M \Coloneqq x & \mid \lambda x: \alpha \mapsto M \mid M \ M \\
    & \mid \texttt{true} \mid \texttt{false} \\
    & \mid \text{if} \ M \ \text{then} \ M \ \text{else} \ M
\end{align*}

Now besides the evaluation rules that we saw before we have to add the 
\emph{typing rules}, which are the rules that define the well-typed expressions 
of the calculus. These rules require a \emph{typing context} normally denoted as
\(\Gamma\), which is a map that associates variables with types. These rules
for example ensure that the if-then-else construct is applied to booleans, and
that the branches of the if-then-else construct have the same type.

\subsubsection{Generalization}

The STLC is the basis of most typed lambda calculi, and can be extended allowing
\emph{terms that depend on types}, this is known as \emph{polymorphism} and
the calculus is known as System F. We can also extend the STLC with
\emph{types that depend on other types}, we will call this later a 
\emph{type constructor}, and the calculus is known as System F\(_\omega\).
If we add \emph{dependent types} we get \(\lambda P\), etc...

We can combine these extensions arbitrarily obtaining different system with 
different expressive power, a taxonomy of these systems is the Barendregt's 
\(\lambda\)-cube. At the apex of generality we have the \emph{Calculus of
Constructions}, which is the basis of Coq and Lean. The following sections will
introduce the concepts of the Calculus of Constructions. For a complete
treatment of the \(\lambda\)-cube we refer the reader to~\cite{lambda_with_types}.

\subsection{Type theory vs set theory}

We follow the type theory developed in~\cite{hottbook}, although any 
Martin-Löf type theory would suffice. In fact getting into the differences 
between Martin-Löf type theory, the Calculus of Constructions, and the
Calculus of Inductive Constructions is beyond the scope of this text, and we
will use the term \enquote{type theory} to refer to all of them.
We can think of the \(\lambda\)-calculus
presented before as a \enquote{syntax} for our type theory, the types will give
meaning to the expressions of the \(\lambda\)-calculus. In fact a long deal
of the next sections will be dedicated to show how to \emph{construct} types
that can encode the object of our proofs.

A key point to remark is that types don't \emph{contain} values, they \emph{classify}
values, while in set theory a value is a \emph{member} of a set, in type theory a
value \emph{has} a type. The analogous to the
membership notation in set theory is the \emph{typing judgement} in type theory,
which is written as \(x: \alpha\), and reads as \enquote{x has type \(\alpha\)}.
While in type theory \(x \in S\) are \(x \notin S\) are both valid expressions,
the typing judgement is \enquote{always true}, it is not a proposition that can be
proven or disproven.
Because of this we cannot introduce a variable without
specifying its type, when we write \(x\) we are implicitly writing \(x: \alpha\),
where \(\alpha\) is the type of \(x\).

The most important difference between type theory and set theory is that type
theory its own deductive system, while set theory is formulated within
first-order logic. This means that set theory has two primitive notions, sets
and propositions, while type theory has types and functions. With just these two
we can recreate all first-order logic, and in fact, we can use types to encode
propositions, and functions to encode proofs. The type judgement \(x: \alpha\)
can also be read as \enquote{x is a proof of \(\alpha\)}, and the function type
\(\alpha \to \beta\) can be read as \enquote{if we have a proof of \(\alpha\) we can
construct a proof of \(\beta\)}. This means that function abstraction corresponds
to a \emph{logical implication} in classical logic, we will expand on this 
in Section~\ref{sec:curry_howard}.

\subsection{Equality}

Equality in type theory is slightly more complex than in classic mathematics,
the equality that we have by default is \emph{definitional equality}. That is two
terms are equal if they can be transformed into each other via \emph{rewriting
rules}. The reduction rules of type theories typically include:

\begin{itemize}
\item \(\beta\)-reduction. Replacing a function application with the function 
  body where the argument is substituted for the bound variable.
\item \(\delta\)-reduction. Replacing a constant with its definition.
\item \(\iota\)-reduction. Simplifying expressions involving constructors and 
  pattern matching for inductive data types.
\end{itemize}

These rules normally satisfy a series of desirable properties, such as
\emph{confluence}, \emph{termination}, and \emph{consistency}, that ensure that 
the reduction process is well defined, and that the equality relation is an 
equivalence relation. In some cases other nice properties are added, such as 
\emph{congruence}, \emph{strong normalization}, and \emph{canonicity}.

Another type of equality is \emph{propositional equality}, which is states that two
terms are equal if they can be proven to be equal in the type theory. This is
more akin to the classic notion of equality, and is used to prove properties of
the terms.\ A special case of propositional equality is \emph{functional
extensionality}, 
which states that two functions are equal if they are
equal for all arguments, a classic result of set theory that follows from the
axiom of extensionality, which states that two sets are equal if they contain
the same elements.\ (The relationship may depend on axiom K).

Another important axiom is \emph{propositional extensionality}, which
states that two propositions are definitionally equal if they are equivalent,
that is, if they imply each other. This is a key property of type theory, as it
allows us to rewrite propositions in proofs, and simplify the reasoning process.

Equality is also a \emph{binary relation}, that is, a relation that relates two
values of a type. In type theory binary relations are functions that take two
values and return a proposition, this proposition is true if the relation holds,
and false otherwise.

\begin{definition}[Relation]
  A relation is a function that takes two elements of a type and returns a
  proposition.
  \[
  \texttt{Rel} \ (\alpha : \texttt{Type}) \ (\beta : \texttt{Type}) 
    \coloneqq \alpha \to \beta \to \Prop 
  \]
\end{definition}

And we can define the relation of equality as:

\begin{definition}[Equality]
  The relation of equality is a function that takes two elements of a type and
  returns a proposition.
  \[
  \texttt{Eq} \ (\alpha: \texttt{Type}) \coloneqq \texttt{Rel} \ \alpha \ \alpha  
  \]
\end{definition}

Equality is much more than an equivalence relation, however. It has the important 
property that every assertion respects the equivalence, in the sense that we can 
substitute equal expressions without changing the truth value. 

\subsubsection{Equivalence}

In type theory, the notion of equivalence is more general than equality, and
refers to a relation between two types that is reflexive, symmetric, and
transitive. This relation is normally defined as a setoid in the type theory,
and is used to define the notion of isomorphism between types.

\subsection{Universes}

Before we introduced the judgement \(x: \alpha\), but now we ask ourselves 
what is the type of \(\alpha\)? The answer is that \(\alpha\) is of type
\(\texttt{Type}\), but what is the type of \(\texttt{Type}\)? Of course it is
also a type so \(\texttt{Type}: \texttt{Type}\), this leads to a problem known as
\emph{impredicativity}.

Impredicativity arises when a definition of an object or type relies on
a quantifier that ranges over a totality encompassing the very object being
defined. This self-referential aspect can lead to logical paradoxes, akin to
Russell's Paradox in set theory, where a set is defined to contain all sets that
do not contain themselves.

The type theory analogous to Russell's paradox is the \emph{Girard's paradox}, in
a system that allows impredicativity, we can define a type \(U\) that contains all
types that do not contain themselves, and then ask if \(U: U\), if it is, then it
should not be of type \(U\), and if it is not, then it should be. A system that
shows this paradox is not very useful for theorem proving, as it is
inconsistent, so any proposition can be proven.\footnote{
  Historically this problem has been solved with System U}. 
To avoid this, we introduce a hierarchy of \emph{type universes}.

Each universe is a \emph{collection of types, whether they are types or not
depends on the particular type theory}, in Lean and Coq they are not. Each
universe contains the types of the previous universe, and is contained in the
next universe. For example: \(\texttt{Type}: \texttt{Type} \ 1\),
\(\texttt{Type} \ 1: \texttt{Type} \ 2\), and so on. The function type inhabits
the smallest universe that contains the types of its domain and codomain. For
example, \(\mathbb{N} \to \mathbb{N}: \texttt{Type} \ 1\), and
\(\texttt{Type} \ 1 \to \texttt{Type} \ 2: \texttt{Type} \ 3\).

The first type universe, \(\texttt{Type} \ 0\), contains the types of the
programming language itself, such as \(\mathbb{N}\), \(\texttt{Bool}\), and
\(\texttt{List}\). The second type universe, \(\texttt{Type} \ 1\), contains the
types of the first universe, and so on.

\subsection{Constructors}

For now we have focused on types, but we need a way to create values of these
types, this is done using \emph{type constructors}, which are functions that
create elements of a type and define the possible forms an element of a type can
take. For example, we can define a type \(\texttt{Bool}\) that represents the
booleans, and define two constructors \(\texttt{true}\) and \(\texttt{false}\) that
create the elements of the type. This ensures that the elements of the type are
either \(\texttt{true}\) or \(\texttt{false}\), and no other value.

\subsection{Pattern matching}

To define functions that operate on types we need a way to distinguish between
the different forms an element of a type can take, this is done using
\emph{pattern matching}. Pattern matching is a way to define functions by
defining a different case for each constructor of the type. For example, we can
define a function \(\texttt{not}\) that takes a boolean and returns its negation
as:
\[
  \texttt{not} \ \texttt{true} = \texttt{false} \qquad
  \texttt{not} \ \texttt{false} = \texttt{true}
\]

\subsection{Inductive types}

Of course defining finite types such as \(\texttt{Bool}\) is not very
interesting, we need a way to define infinite types, such as the natural
numbers. This is done using \emph{inductive types}, which are types that are defined
by a set of constructors that define the possible forms an element of the type
can take. For example, we can define the natural numbers as an inductive type
with two constructors, \(0\) and \(\texttt{succ}\), that create the
elements of the type, in a way analogous to the \emph{Peano axioms}.
\[
0 : \mathbb{N} \qquad
    \texttt{succ}: \mathbb{N} \to \mathbb{N}
\]
Similar to the classic principle of induction, we have defined a base case
\(0\) and an inductive case \(\texttt{succ}\), that creates the
successor of a natural number. In this way we can create the natural numbers as
an inductive type, and define functions that operate on them, such as addition,
multiplication, and exponentiation. This is the essence of inductive types, the
ability to define types that are created by a set of constructors, and define
functions that operate on them. And as is the case in classical mathematics, we
can use induction to prove properties of these types, in this case \emph{structural
induction}.

Following the natural number example, let's construct a function
\(f: \mathbb{N} \to C\) by recursion on the naturals, it's enough to provide a
base case \(c_0 : C\) and a \enquote{next step} function \(c_s: \mathbb{N} \to C \to C\),
then we can define \(f\) via pattern matching as:
\begin{align*}
  f \ 0 & = c_0 \\
  f \ (\texttt{succ} \ n) & = c_s \ n \ (f \ n)
\end{align*}
To use this proof technique we must ensure that the inductive type is
\emph{well-founded}, that is, that the constructors of the type are applied a finite
number of times to create an element of the type.

\subsection{Dependent function type}

In type theory, functions are first class citizens, and can be passed as
arguments, returned as results, and stored in data structures. In some cases we
need to define functions whose return type depends on the value of the
arguments, this is done using \emph{Pi types}, which are types that depend on a
value.

The notation for a Pi type is \(\Pi_{x: \alpha} B \ x\), where \(x\) is the argument
on which the type depends, \(A\) is the type of the argument, and \(B\) is a family
of types indexed by \(x\) (\(B: \alpha \to \texttt{Type}\)). For example, we can
define a function that takes a natural

To better understand this we introduce an example were \(\Pi\)-types are used to
to define a vector of a given length, which type depends on the length of the
vector.

\begin{example}[Vector]
  We can define a vector of a given length as an inductive type that depends on
  the length of the vector. We define the type \(\texttt{Vec}\) as:
\[
 \Pi_{n: \mathbb{N}} \texttt{Vec} \ \alpha \ n
\]
with constructors:
\begin{align*}
  \texttt{nil} & : \texttt{Vec} \ \alpha \ 0 \\
  \texttt{cons} & : \Pi_{n: \mathbb{N}}
  (\alpha \to \texttt{Vec} \ \alpha \ n
    \to \texttt{Vec} \ \alpha \ (\texttt{succ} \ n)) 
\end{align*}
\end{example}
In this way \(\Pi\)-types give us a way to type complex data structures, such as
lists, trees, and graphs, that depend on the values of their elements,
ensuring type safety.

\subsection{Product types}

\subsection{Sum types}

\subsection{The Curry-Howard isomorphism}\label{sec:curry_howard}

The Curry-Howard isomorphism is a result that establishes a correspondence
between \emph{propositions} and \emph{types} and between \emph{terms} and \emph{proofs}.
This is know as the \emph{propositions as types} principle, and is the basis of
all modern proof assistants, such as Coq, Lean, and Agda. In our type system
a proposition is a type \(\Prop\), \enquote{proofs} are terms of that type,
and the proposition is true if the type is inhabited, that is, if there is a
term of that type, this is our \emph{witness}.

Proofs of terms can be manipulated in classical logic, for example we can
obtain a proof that two propositions hold by using the \(\land\) operator,
which states that if \(P\) and \(Q\) are true, then \(P \land Q\) is true. In
type theory this corresponds to the \emph{pair type} \(P \times Q\), which
states that if we have a proof of \(P\) and a proof of \(Q\), then we can
construct a proof of \(P \times Q\). The \(\bot\) judgement which 
represent unprovability is represented as a type with no constructors,
and the \(\top\) judgement which represents a tautology is represented as the unit 
type, which has only one constructor with no arguments.

Let's consider the example of implication (\(\implies\)) in propositional
logic, the proposition \(P \implies Q\) states that if \(P\) is true, then \(Q\)
must be true. In type theory this corresponds to the function type
\(P \to \ Q\), which states that if we have a proof of \(P\), then
we can construct a proof of \(Q\). The proof of the proposition corresponds to a
term of the type, and the proof of the implication corresponds to a function
that takes a proof of \(P\) and returns a proof of \(Q\).

This correspondence is not limited to propositional logic, but extends to
predicate logic, and higher order logics. For example, the universal quantifier
\(\forall\) corresponds to the dependent function type \(\Pi\), and the existential
quantifier \(\exists\) corresponds to the dependent pair type \(\Sigma\). This
correspondence is the basis of all modern proof assistants, and allows us to
reason about programs using the tools of logic, and vice versa.

The \emph{type checker} of the system then allows us to \emph{verify} the proofs, by
checking that the terms are well-typed, and that the terms correspond to valid
proofs. This mechanism ensures that the proofs are correct, and that the
programs are well-typed.

\subsection{Sets}\label{sec:sets}

We can express a set in type theory as a function that takes an element of the
type and returns a proposition, this propsition tells us if the element is in
the set or not.

\begin{definition}[Set]
  A set is a function that takes an element of a type and returns a proposition.
  \[
  \texttt{Set} \ (\alpha: \texttt{Type}) \coloneqq \alpha \to \Prop 
  \]
\end{definition}

We can also use the familiar \emph{set builder} notation to define sets.
\begin{notation}[Set builder]
  We can define a set using the set builder notation as:
  \[
  \{x \mid P\} \coloneqq \lambda x \mapsto P 
  \]
\end{notation}

The definition of the membership relation is immediate as the definition of the
set itself captures the concept of membership.

\begin{definition}[Membership]
  The membership relation is a function that takes an element of a type and a set
  and returns a proposition, let \(x: \alpha\) and \(S: \texttt{Set} \ \alpha\):
  \[
  x \in S \coloneqq S \ x 
  \]
\end{definition}

A subset is defined in a way similar to classic set theory.

\begin{definition}[Subset]
  A subset is a function that takes an element of a type and two sets and returns
  a proposition, let \(A \ B: \texttt{Set} \ \alpha\):
  \[
  A \subseteq B \coloneqq \forall x, x \in A \implies x \in B 
  \]
\end{definition}

And with this a strict subset is defined as:

\begin{definition}[Strict subset]
  A strict subset is subset different from the set itself, let
  \(A \ B: \texttt{Set} \ \alpha\), then:
  \[
  A \subset B \coloneqq A \subseteq B \land A \neq B
  \]
\end{definition}

With this definitions set operations are defined in exactly the same way as in
classic set theory, and we can reason about sets using the tools of type theory.

\begin{definition}[Set operations]
  We can define the set operations as:
  \begin{align*}
    A \cup B & \coloneqq \{x \mid x \in A \lor x \in B\} \\
    A \cap B & \coloneqq \{x \mid x \in A \land x \in B\} \\
    A \setminus B & \coloneqq \{x \mid x \in A \land x \notin B\} \\
    A^\complement & \coloneqq \{x \mid x \notin A\} 
  \end{align*}
\end{definition}

So we can also use sets to reason if necessary, and it will be, as we will see
in the definition of denotational semantics.

\section{Proof irrelevance}

We've seen that Lean and Coq share much of their type theory, so what are the
key differences between the two? The most important difference is that Lean is
designed to be proof-irrelevant. Proof-irrelevance is a property of a type
theory that states that all proofs of a proposition are definitionally equal,
meaning that they are indistinguishable by the type checker.

For example, imagine we have two proofs of the fact that \(2 + 2 = 4\), one that
uses the Peano axioms, and another that uses the definition of addition. In a
proof-irrelevant type theory, these two proofs would be considered equal, and
could be used interchangeably.

Why does this matter? First of all it helps to simplify complex proofs, as we
don't have to worry about the details of the proof, only that it exists. This
also has computational benefits, as we can erase the proofs from the final
program, and only keep the results.

It also aids in \emph{code extraction}, the process of extracting a program from
proofs and definitions. This is a key feature when using Lean as tool in formal
verification, as we can first formally prove the correctness of the program, and
then extract an efficient implementation from the proof. Irrelevancy ensures
that the extracted code is independent of the proof used. Code extraction is
an area where Lean 4 has made significant improvements over Lean 3, as it
supports a better compiler backend, and more target languages.

Proof irrelevance is also supported in Coq, but not as a definitional equality,
but a set of axioms that ensure that all proofs of a proposition are equal. In
Coq \(\Prop\) is a universe, separate from data types which reside in
\(\texttt{Type}\), inside this universe two proofs of the same proposition are
\emph{propositionally equal}. Coq also supports definitional proof irrelevance, by
using the \(\texttt{SProp}\) but it's not the default.

\section{Big-Step Semantics}

Now that we have introduced type theory and proof assistants, we can move on to
using them to define the semantics of programming languages. The exposition
will follow closely the one in~\cite{winskel_formal_semantics}.
Semantics are the
study of the \emph{meaning} of programs, and can be defined in different ways, we
will focus on operational semantics, and denotational semantics.

We first describe a simple imperative language \(\texttt{IMP}\) that we will use
to illustrate the concepts.\ \(\texttt{IMP}\) is a simple language that consists of
a set of commands, such as assignment, sequencing, conditionals, and loops. It
is constrained by design to be simple, so that we can focus on introducing the
techniques of formal semantics, but this approach can be extended to widely used
languages, such as C, Java, and Python.

\subsection{Syntax}

As an informal introduction to the syntax of \(\texttt{IMP}\) we can define the
\enquote{syntactic sets} of the language as:

\begin{itemize}
\item Integers \(n \in \mathbb{Z}\).
\item Boolean values \(b \in \texttt{Bool} = \{\texttt{true}, \texttt{false}\}\).
\item Variables \(x \in \texttt{Var}\), that store integers.
\item Arithmetic expressions \(a \in \texttt{Aexp}\), that can be constants,
  variables, or arithmetic operations.
\item Boolean expressions \(b \in \texttt{Bexp}\), that can be constants, variables,
  or boolean operations.
\item Commands \(c \in \texttt{Com}\), that can be assignments, conditionals, loops,
  or sequences of commands.
\end{itemize}

Now we could define the syntax of \(\texttt{Aexp}\) in Backus-Naur form:
\[
a \Coloneqq n \ | \ x \ | \ a_1 + a_2 \ | \ a_1 - a_2 \ | \ a_1 \times a_2
\]
the syntax of \(\texttt{Bexp}\) as:
\[
b \Coloneqq \texttt{true} \ | \ \texttt{false} \ | \ 
  !b \ | \ b_1 \land b_2 \ | \ b_1 \lor b_2 \ | \ a_1 = a_2 \ | \ a_1 \leq a_2
\]
and the syntax of \(\texttt{Com}\) as:
\[
c \Coloneqq \texttt{skip} \ | \ x \ \texttt{:=} \ a \ | \ c_1; c_2 \ | \
      \texttt{if} \ b \ \texttt{then} \ c_1 \ \texttt{else} \ c_2 \ \texttt{end} 
      \ | \ \texttt{while} \ b \ \texttt{loop} \ c \ \texttt{end}
\]
In type theory we can encode this as inductive types, the type of arithmetic 
expressions can be defined as an inductive type \(\texttt{Aexp}\) with 
constructors for each \enquote{syntactic construct}:
\begin{align*}
  \texttt{val} &\coloneqq \mathbb{Z} \to \texttt{Aexp} \\
  \texttt{var} &\coloneqq \texttt{String} \to \texttt{Aexp} \\
  \texttt{add} &\coloneqq \texttt{Aexp} \to \texttt{Aexp} \to \texttt{Aexp} \\
  \texttt{sub} &\coloneqq \texttt{Aexp} \to \texttt{Aexp} \to \texttt{Aexp} \\
  \texttt{mul} &\coloneqq \texttt{Aexp} \to \texttt{Aexp} \to \texttt{Aexp}
\end{align*}
While the first to simply \enquote{coerce} a value or a variable into an
expression, the others \enquote{construct} an expression from two other
expressions. The constructors for boolean expressions are defined in a similar
fashion:
\begin{align*}
  \texttt{true} &\coloneqq \texttt{Bexp} \\
  \texttt{false} &\coloneqq \texttt{Bexp} \\
  \texttt{not} &\coloneqq \texttt{Bexp} \to \texttt{Bexp} \\
  \texttt{and} &\coloneqq \texttt{Bexp} \to \texttt{Bexp} \to \texttt{Bexp} \\
  \texttt{or} &\coloneqq \texttt{Bexp} \to \texttt{Bexp} \to \texttt{Bexp} \\
  \texttt{eq} &\coloneqq \texttt{Aexp} \to \texttt{Aexp} \to \texttt{Bexp} \\
  \texttt{le} &\coloneqq \texttt{Aexp} \to \texttt{Aexp} \to \texttt{Bexp}
\end{align*}
And the same goes for commands:
\begin{align*}
  \texttt{skip} &\coloneqq \texttt{Com} \\
  \texttt{assign} &\coloneqq \texttt{String} \to \texttt{Aexp} \to \texttt{Com} \\
  \texttt{seq} &\coloneqq \texttt{Com} \to \texttt{Com} \to \texttt{Com} \\
  \texttt{ifElse} &\coloneqq \texttt{Bexp} \to \texttt{Com} \to \texttt{Com} \to \texttt{Com} \\
  \texttt{whileLoop} &\coloneqq \texttt{Bexp} \to \texttt{Com} \to \texttt{Com}
\end{align*}

Now a simple expression in our language would look something like:
\[
  \texttt{add} \ (\texttt{var} \ \textit{``x''}) \ (\texttt{val} \ 3)
\]
this is not very readable, so we will use some \enquote{abuse of notation} and
say that \(x + 3\) is an arithmetic expression of \texttt{IMP} for example,
when defining our language in Lean we have to be more careful about mixing
our language and meta-language, by using some facilities provided by Lean's 
macros we can make a little parser for our language, and by using type
classes and coercions we can even use some familiar notation in the host language,
all of the details can be found in 
\url{https://github.com/remind-me-later/SemanticsLean/blob/main/Semantics/Imp/Lang.lean}.

An example program in \texttt{IMP} could then be:

\begin{example}[Factorial in \texttt{IMP}]
  The factorial of a number can be computed using a simple loop, we can define
  the factorial of a number \(n\) as:
  \begin{verbatim}
  n := 5; i := 2; res := 1;

  while i <= n loop
      res := res * i;
      i := i + 1
  end
  \end{verbatim}
\end{example}

\subsection{States}

As we've seen in the syntax our imperative language supports a feature that
characterizes all imperative languages are states, and assignments, the
definition of states is more nuanced that one could expect at first. A state is
a function that maps variables to values, we could define it as:
\[
\funfont{state} : \tyfont{String} \to \mathbb{Z}
\]
But we will take a more roundabout way, and define first a \(\tyfont{Map}\)
type.

\begin{definition}[Map]
  A map is a function that maps keys (strings) to values, we can define it as:
  \[
  \tyfont{Map} \ (\alpha : \tyfont{Type}) \coloneqq \tyfont{String} \to \alpha
  \]
\end{definition}

Of course this map is not particularly useful if we don't have a way to relate
keys to values, this is done using the \texttt{update} function:

\begin{definition}[Map update]
  The update function takes a map, a key, and a value, and returns a new map
  that maps the key to the value, and all other keys to the same value as the
  original map.
  \begin{align*}
  \funfont{update} \ & (m: \tyfont{Map} \ \alpha) \ (k: \tyfont{String}) \ (v: \alpha): \tyfont{Map} \ \alpha \coloneqq \\
                      &\lambda k' \mapsto \funfont{cond} \ (k = k') \ v \ (m \ k')
  \end{align*}
\end{definition}

The \funfont{cond} function is nothing more than syntactic sugar for 
pattern matching, we can define it as:
\begin{align*}
  \funfont{cond} \ \texttt{true} \ v \ w &\coloneqq v \\
  \funfont{cond} \ \texttt{false} \ v \ w &\coloneqq w
\end{align*}
We can even denote the \funfont{cond} function as an \emph{if then else}:
\begin{align*}
  \funfont{update} \ & (m: \tyfont{Map} \ \alpha) \ (k: \tyfont{String}) \ (v: \alpha): \tyfont{Map} \ \alpha \coloneqq \\
                      &\lambda k' \mapsto \text{if} \ k = k' \ \text{then} \ v \ \text{else} \ m \ k'
\end{align*}
Not to be confused with the \texttt{if then else} construct of the language 
\texttt{IMP}. While we are introducing notation we can also introduce some to 
denote the \enquote{assignments} of a map in a more familiar way:
\begin{notation}[Map update]
  We will denote \((\funfont{update} \ m \ k \ v)\) as \(m[k \leftarrow v]\).
\end{notation}

With these definitions we can prove a lot of interesting properties that hold
in the majority of programming languages, for example:

\begin{theorem}
  If we assign a value to a variable, the map evaluates to that value.
  \[
  (m[k \leftarrow v]) \ k = v
  \]
\end{theorem}

\begin{theorem}
  If we assign a value to a variable, the map does not evaluate to that value
  for other variables (\(k \neq k'\)).
  \[
  (m[k \leftarrow v]) \ k' = m \ k'
  \]
\end{theorem}

Both these fact are pretty intuitive, let's see another theorem that is also
intuitive but not so easy to prove. Usually assigning a value to a
variable makes the computer \enquote{forget} the previous value. With the
\funfont{update} function as defined above this is not true by definition,
but we can easily prove it.

\begin{theorem}
  If we assign a value to a variable, and then assign another value to the same
  variable, the map \enquote{forgets} the first value.
  \[
  m[k \leftarrow v_2][k \leftarrow v_1] = m[k \leftarrow v_1]
  \]
\end{theorem}

\begin{proof}
  To prove the equality of both functions we use the axiom of \emph{function
  extensionality}, this transforms the problem into:
  \[
  \forall k' , (m[k \leftarrow v_2][k \leftarrow v_1]) \ k' = (m[k \leftarrow v_1]) \ k'
  \]
  Now we can distinguish two cases:
  \begin{itemize}
    \item[(\(k = k'\))] In this case the equality holds by \enquote{unfolding} the
      definition of the \funfont{update} function.
      \[
        \text{if} \ k = k' \ \text{then} \ v_1 \ \text{else} \ 
          \text{if} \ k = k' \ \text{then} \ v_2 \ \text{else} \ m \ k' = 
          \text{if} \ k = k' \ \text{then} \ v_1 \ \text{else} \ m \ k'
      \]
      noting that \(k = k'\) is true, we can simplify the expression to \(v_1 = v_1\).
    \item[(\(k \neq k'\))] This case is analogous.
  \end{itemize}
\end{proof}

The proof of other similar properties, such as the commutativity of the update
function, or the idempotence of the update function, are proved in a similar
fashion, and can be found in the code.

\begin{theorem}
  If we assign a value to a variable, and then assign another value to another
  variable (\(k_1 \neq k_2\)), the order of the assignments does not matter.
  \[
  m[k_2 \leftarrow v_2][k_1 \leftarrow v_1] = m[k_1 \leftarrow v_1][k_2 \leftarrow v_2]
  \]
\end{theorem}

\begin{theorem}
  If we assign to a key the value it already has, the map does not change.
  \[
  m[k \leftarrow m \ k] = m
  \]
\end{theorem}

With these properties proven we can finally define our \(\tyfont{State}\).

\begin{definition}[State]
  A state is a map that maps variables to values, we can define it as:
  \[
  \tyfont{State} \coloneqq \tyfont{Map} \ \mathbb{Z}
  \]
\end{definition}

Now we have a problem, what value do the variables have when the program starts?
In some popular programming languages, like C, wen variables are read before
they are assigned the hold a random value, this value depends on whatever
the contents of the memory now used to store the variable used to hold, we could
model this by defining an \emph{undefined} state, but the \(\tyfont{Map}\)
that we have defined is \emph{total}, that is, it is defined for all inputs. 
We could define a partial map:

\begin{definition}[Partial Map]
  A partial map is a map that is not defined for all inputs, we can define
  it as:
  \[
  \tyfont{PMap} \ (\alpha : \tyfont{Type}) \coloneqq \tyfont{String} \to \tyfont{Option} \ \alpha
  \]
\end{definition}

But this would complicate the definitions, and proofs, needlessly so we will
take the simpler approach of setting all variables to zero in the \emph{initial state},
of course, setting them to zero is just a convention, we could set them to any
value.

\begin{definition}[Initial state]
  We will define the \emph{initial state} of an \(\texttt{IMP}\) program as the
  constant function that maps all variables to zero.
  \[
  s_0 \coloneqq \lambda x \mapsto 0
  \]
\end{definition}

With all this evaluating the state of a program is as simple as applying it to
a string:
\[
(s_0[k \leftarrow 3][k \leftarrow 4][k \leftarrow 7]) \ k = 7
\]

\subsection{Evaluating expressions}

Now that we have in place the syntax of the language, we can start to wonder
what does it mean to evaluate an expression. Because our expressions
contain variables, they don't have any meaning without a state associated to
them, so our evaluation will act on a pair of an expression and a state 
\(\tyfont{Aexp} \times \tyfont{State}\), we 
will call this pair a \emph{configuration}.

Our evaluation will be \emph{syntax-directed}, that is, for every syntactic
construct of the language we will define a rule that tells us how to evaluate
it. This rule will be a \emph{binary relation} between configurations and values,
we will call this relation \(\funfont{big\_step}\), and define it as:

\begin{definition}[Step relation]
  The step relation is a binary relation between configurations and values, it 
  has type:
  \[
  \funfont{big\_step} : \tyfont{Aexp} \times \tyfont{State} \to \mathbb{Z} \to \Prop
  \]  
\end{definition}

This is what is generally known as \emph{big step semantics, as it evaluates
the expression to a value in a single step, from configuration directly to value.}
Let's introduce a notation more familiar to classic text on formal semantics.
\begin{notation}[Arithmetic big step]
  We will denote the step relation \(\funfont{big\_step} \ (a, s) \ n\) as
  \[
  (a, s) \bigstep n
  \]
\end{notation}

Now the constructors will be the rules of the evaluation, for example the
constructor for the constant expression will be:
\[
  \texttt{val} : (n, \cdot) \bigstep n
\]

Which relates the configuration \((n, \cdot)\), where the \(n\) of the left is
a constant expression (in reality we should write \(\texttt{val} n\) 
to denote that its the first constructor of the \(\tyfont{Aexp}\) type, but
we will omit it for brevity) and the \(\cdot\) denotes any state, to the
integer value \(n\). Some attention is needed to not confuse the constructor
of the \emph{syntactic expression} with that of the \emph{evaluation relation},
but the distinction is clear from the context.

The constructor for the variable expression is similar, but this time we need
to look up the value of the variable in the state:
\[
  \texttt{var} : (x, s) \bigstep s \ x
\]

Now we can define the constructors for the arithmetic operations, since these 
are \emph{binary} operations we will need to know the value of the left and
right expressions to evaluate the operation, this is done by requiring
the evaluation of the left and right expressions in the premises of the rule.
For example the constructor for the addition operation is:
\[
  \texttt{add} : \frac{(a_1, s) \bigstep n_1 \quad (a_2, s) \bigstep n_2}
   {(a_1 + a_2, s) \bigstep n_1 + n_2}
\]
The previous rule states that if the configuration \((a_1, s)\) evaluates to
\(n_1\) and the configuration \((a_2, s)\) evaluates to \(n_2\), then the
configuration \((a_1 + a_2, s)\) evaluates to \(n_1 + n_2\).

The constructors of the subtraction and multiplication operations are analogous.
\begin{align*}
  \texttt{sub} &: \frac{(a_1, s) \bigstep n_1 \quad (a_2, s) \bigstep n_2}
   {(a_1 - a_2, s) \bigstep n_1 - n_2} \\
  \texttt{mul} &: \frac{(a_1, s) \bigstep n_1 \quad (a_2, s) \bigstep n_2}
   {(a_1 * a_2, s) \bigstep n_1 * n_2}
\end{align*}

One can think that while applying the syntactic rules of the language we are
\enquote{building up} the expression while we apply the evaluation rules
we will \enquote{break down} the expression to its value.
Let's see an example of a valid \emph{evaluation tree}.

\begin{example}[Arithmetic evaluation]\label{ex:arithmetic_evaluation}
  Let's evaluate the value of the expression
  \((x + 5) + (9 - 7)\) in the state \(s_0\) by applying the \enquote{rules} 
  (constructors) of the evaluation relation:
\[
  \texttt{add}: \frac{
    \texttt{add}: \displaystyle{\frac{
      \texttt{var}: (x, s_0) \bigstep 0
      \quad
      \texttt{val}: (5, s_0) \bigstep 5
    }{(x + 5, s_0) \bigstep 5}}
    \quad 
    \texttt{sub}: \displaystyle{\frac{
      \texttt{val}: (9, s_0) \bigstep 9
      \quad
      \texttt{val}: (7, s_0) \bigstep 7
    }{(9 - 7, s_0) \bigstep 2}}
  }
  {((x + 5) + (9 - 7), s_0) \bigstep 7}
\]
\end{example}

Another way to express evaluation is via a \emph{recursive function} that breaks
down the expression into its components and evaluates them, the rules are
analogous to the ones we have defined before, but this time the function will
take a configuration and return a value.

\begin{definition}[Arithmetic evaluation]
  The evaluation function takes a configuration and returns a value, it has type:
  \[
  \funfont{eval} : \tyfont{Aexp} \to \tyfont{State} \to \mathbb{Z}
  \]
\end{definition}

This time, instead of a Cartesian product, we will use a \emph{curried} function.
We can easily define the evaluation rules with pattern matching:
\begin{align*}
  \funfont{eval} \ (\texttt{val} \ v) \ \cdot &\coloneqq v \\
  \funfont{eval} \ (\texttt{var} \ x) \ s &\coloneqq s \ x \\
  \funfont{eval} \ (\texttt{add} \ a_1 \ a_2) \ s &\coloneqq \funfont{eval} \ a_1 \ s + \funfont{eval} \ a_2 \ s \\
  \funfont{eval} \ (\texttt{sub} \ a_1 \ a_2) \ s &\coloneqq \funfont{eval} \ a_1 \ s - \funfont{eval} \ a_2 \ s \\
  \funfont{eval} \ (\texttt{mul} \ a_1 \ a_2) \ s &\coloneqq \funfont{eval} \ a_1 \ s * \funfont{eval} \ a_2 \ s
\end{align*}

We can introduce some notation to denote the evaluation of an expression on a state.

\begin{notation}[Arithmetic evaluation]
  We will denote the evaluation function (\(\funfont{eval} \ a \ s\)) as applying
  the expression \(a\) to the state \(s\):
  \[
  a \ s
  \]
\end{notation}

Now the example~\ref{ex:arithmetic_evaluation} can the be directly 
evaluated by the \(\funfont{eval}\) function:
\[
  ((x + 5) + (9 - 7)) \ s_0 = 7
\]

This approach has some clear advantages over the relation definition, as
it can compute the result of the evaluation directly, instead of 
just saying if some \enquote{evaluation tree} is valid or not. That is only if
both evaluation rules are equivalent, fortunately they are, and we can prove it.

\begin{theorem}
  The step relation and the \(\funfont{eval}\) function are equivalent, that is, for all
  expressions \(a\), states \(s\), and values \(n\), we have:
  \[
  ((a, s) \bigstep n) = (a \ s = n)
  \]
\end{theorem}

\begin{proof}
The proof, follows by applying \emph{propositional extensionality} to
transform the equals into a double implication, and then proving each
implication separately.

Let's first prove the left to right implication 
\((a, s) \bigstep n \implies a \ s = n\). We proceed by induction
on the evaluation relation, the cases for the constructors \(\texttt{val}\)
and \(\texttt{var}\) are immediate. In the case of \(\texttt{add}\) the
expression has the form \(a_1 + a_2\), and so we have to prove that
\[
(a_1 + a_2 \ s) = n_1 + n_2
\]
for some \(n_1\) and \(n_2\), by the induction hypothesis we have that
\(a_1 \ s = n_1\) and \(a_2 \ s = n_2\), and so we just have to apply the
definition of the \(\funfont{eval}\) function to prove the equality.
The cases for the subtraction and multiplication are analogous.

The right to left implication \(a \ s = n \implies (a, s) \bigstep n\)
is proven by induction on arithmetic expression, again the cases for 
\(\texttt{val}\) and \(\texttt{var}\) are immediate, so let's focus on
\(\texttt{add}\). Like before we have to prove that 
\((a_1 + a_2, s) \bigstep n\), with the induction hypotheses:
\[
\begin{gathered}
  \forall n, a_1 \ s = n \implies (a_1, s) \bigstep n \qquad
  \forall n, a_2 \ s = n \implies (a_2, s) \bigstep n
\end{gathered}
\]
by choosing \(n\) to be \(a_1 \ s\) and \(a_2 \ s\) respectively we can construct
the evaluation tree:
\[
  \texttt{add}: \frac{
    (a_1, s) \bigstep a_1 \ s
    \quad
    (a_2, s) \bigstep a_2 \ s
  }
  {
   (a_1 + a_2, s) \bigstep a_1 \ s + a_2 \ s
  }
\]
And now we simply apply the hypothesis (\((a_1 + a_2) \ s = n\)) to obtain the result. The cases for
subtraction and multiplication are analogous.
\end{proof}

We present the rules of evaluation of boolean expression without much detail,
as they are analogous to the arithmetic expressions, we have presented first
the rules for arithmetic expressions because some of the rules for boolean
expressions depend on the evaluation of arithmetic expressions.
\[
\begin{gathered}
  \texttt{true}: (\texttt{true}, \cdot) \bigstep \texttt{true} \quad
  \texttt{false}: (\texttt{false}, \cdot) \bigstep \texttt{false} \\
  \texttt{not}: \frac{(b, s) \bigstep b'}{\lnot (b, s) \bigstep \lnot b'} \quad
  \texttt{and}: \frac{(b_1, s) \bigstep b_1' \quad (b_2, s) \bigstep b_2'}
    {(b_1 \land b_2, s) \bigstep b_1' \land b_2'} \quad
  \texttt{or}: \frac{(b_1, s) \bigstep b_1' \quad (b_2, s) \bigstep b_2'}
    {(b_1 \lor b_2, s) \bigstep b_1' \lor b_2'} \\
  \texttt{eq}: (a_1 = a_2, s) \bigstep a_1 \ s = a_2 \ s \quad
  \texttt{le}: (a_1 \leq a_2, s) \bigstep a_1 \ s \leq a_2 \ s
\end{gathered}
\]

Notice that in the last two rules we are using the (\funfont{eval}) function to evaluate
the arithmetic expressions, we could have used the evaluation relation, but
this approach is easier. Everything that we said about the (\funfont{eval}) function
and its equivalence to the big step relation holds for boolean expressions.

\subsection{Evaluating commands}

In the previous section defined the big step semantics for arithmetic and
boolean expressions, and also their reduction functions, one could wonder
why do we bother with introducing the big step semantics since 
the functions are easier to define and use. The reason is that 
\emph{commands} cannot be evaluated by a function, this is due to the 
\texttt{while} construct. We could define a \(\funfont{eval}: \tyfont{Com} \to \tyfont{State} \to \tyfont{State}\) 
function to evaluate commands that looks something like this:
\begin{align*}
  \funfont{eval} \ \texttt{skip} \ s &\coloneqq s \\
  \funfont{eval} \ (x \ \texttt{:=} \ a) \ s &\coloneqq s[x \leftarrow a \ s] \\
  \funfont{eval} \ (c_1; c_2) \ s &\coloneqq \funfont{eval} \ c_2 \ (\funfont{eval} \ c_1 \ s) \\
  \funfont{eval} \ (\texttt{if} \ b \ \texttt{then} \ c_1 \ \texttt{else} \ c_2 \ \texttt{end}) \ s &\coloneqq
    \text{if} \ b \ s \ \text{then} \ \funfont{eval} \ c_1 \ s \ \text{else} \ \funfont{eval} \ c_2 \ s \\
  \funfont{eval} \ (\texttt{while} \ b \ \texttt{loop} \ c \ \texttt{end}) \ s &\coloneqq
    \text{if} \ b \ s \ \text{then} \ \funfont{eval} \ (\texttt{while} \ b \ \texttt{loop} \ c \ \texttt{end}) \ (\funfont{eval} \ c \ s) \ \text{else} \ s
\end{align*}
The problem with this function is that it is not \emph{well founded}, that is,
it could loop indefinitely, for example, the evaluation of the command
\texttt{while true loop skip end}
would never terminate. This is why we
need a relation to define the semantics of commands, later we will see
an approach that will give us a \enquote{function} to evaluate commands, that
is \emph{denotational semantics}.

To define the big step semantics of commands we will use a relation, as 
before, this time associating configurations to states. The definition of
the constructors for skip, assignment and concatenation are straightforward:
\begin{align*}
  \texttt{skip} &: (\texttt{skip}, s) \bigstep s \\
  \texttt{ass} &: (v \ \texttt{:=} \ a, s) \bigstep s[v \leftarrow a \ s] \\
  \texttt{cat} &: \frac{
    (c_1, s_1) \bigstep s_2 \quad (c_2, s_2) \bigstep s_3
  }{
    (c_1; c_2, s_1) \bigstep s_3
  }
\end{align*}
In the concatenation we have to introduce an intermediate state \(s_2\) that doesn't appear in the
final result, in this way we account for \enquote{intermediate step} in the 
evaluation.

The rules for the \texttt{if} construct require two constructors,
one for the case when the condition is false:
\[
\texttt{ifFalse} : \frac{b \ s_1 = \texttt{false}  \quad (c_1, s_1) \bigstep s_2}
{(\texttt{if} \ b \ \texttt{then} \ c_1 \ \texttt{else} \ c_2 \ \texttt{end}, s_1) \bigstep s_2}
\]
and another for the case when the condition is true:
\[
\texttt{ifTrue} : \frac{b \ s_1 = \texttt{true}  \quad (c_2, s_1) \bigstep s_2}
{(\texttt{if} \ b \ \texttt{then} \ c_1 \ \texttt{else} \ c_2 \ \texttt{end}, s_1) \bigstep s_2}
\]
The evaluation rule for \texttt{while} follows the same principle but with a more
recursive approach, the case when the condition is false is easy:
\[
\texttt{whileFalse} : \frac{b \ s = \texttt{false}}
  {(\texttt{while} \ b \ \texttt{loop} \ c \ \texttt{end}, s) \bigstep s}
\]
but when the condition is true, we have to recurse from another evaluation tree
where the \enquote{previous step} has already been applied:
\[
\texttt{whileTrue} : \frac{b \ s_1 = \texttt{true} \quad (c, s_1) \bigstep s_2 \quad
  (\texttt{while} \ b \ \texttt{loop} \ c \ \texttt{end}, s_2) \bigstep s_3}
{(\texttt{while} \ b \ \texttt{loop} \ c \ \texttt{end}, s_1) \bigstep s_3}
\]
While the evaluation rules for
expressions could be applied \enquote{automatically}, the evaluation rules for
commands require a more \enquote{manual} approach, as we have to \enquote{choose}
the correct constructor to apply in the control flow of the program.

Introducing some notation, as per usual, we can check the evaluation of a
program like this:
\begin{example}[Program evaluation]
  Let's evaluate a simple program containing a branch in the state \(s_0\):
\begin{verbatim}
  x := 2;
  if x <= 1 then
    y := 3
  else
    z := 4
  end
\end{verbatim}
  to simplify the notation we will denote the program as \(P\), and the if construct as \(\text{Ifelse}\), then:
  \[
  \texttt{cat}: \frac{
    \texttt{ass}: (\texttt{x := 2}, s_0) \bigstep s_0[x \leftarrow 2]
    \quad
    \texttt{ifFalse}: \displaystyle{\frac{
      \texttt{ass}: (\texttt{z := 4}, s_0[x \leftarrow 2]) \bigstep s_0[x \leftarrow 2][z \leftarrow 4]
    }
    {
      (\text{Ifelse}, s_0[x \leftarrow 2]) \bigstep s_0[x \leftarrow 2][z \leftarrow 4]
    }}
  }{
    (P, s_0) \bigstep s_0[x \leftarrow 2][z \leftarrow 4]
  }
  \]
\end{example}
As we can see the \enquote{proof} that the program finishes in the given state
follows directly from the evaluation tree created by the constructors of the
big step semantics.

\subsection{Rewriting rules}

Now we can prove some \enquote{rewriting rules} for the constructs of the 
language, the proof is more or less immediate and can be found in the code,
so we will only present the theorems.

\begin{theorem}
  Skipping (doing nothing) in a program doesn't change the state:
  \[
  ((\texttt{skip}, s_1) \bigstep s_2) = (s_1 = s_2)
  \]
\end{theorem}

\begin{theorem}
  The concatenation of two programs is equivalent to executing the first
  program and then the second:
  \[
  ((c_1 \texttt{;} c_2, s_1) \bigstep s_3) = 
  \exists s_2, \left((c_1, s_1) \bigstep s_2 \land (c_2, s_2) \bigstep s_3\right)
  \]
\end{theorem}

The rewriting rules for the \texttt{if} construct can be defined in two equivalent ways,
which rule to apply depends on the context, as either one could be easier to use
depending on the situation.

\begin{theorem}\label{thm:if_rewrite_1}
  The \texttt{if} construct can be rewritten as:
  \[
  ((\texttt{if} \ b \ \texttt{then} \ c_1 \ \texttt{else} \ c_2 \ \texttt{end}, s_1) \bigstep s_2) = 
  \text{if} \ b \ s_1 \ \text{then} \ (c_1, s_1) \bigstep s_2 \ \text{else} \ (c_2, s_1) \bigstep s_2
  \]
\end{theorem}
or, equivalently:
\begin{theorem}\label{thm:if_rewrite_2}
  The \texttt{if} construct can be rewritten as:
  \[
  ((\texttt{if} \ b \ \texttt{then} \ c_1 \ \texttt{else} \ c_2 \ \texttt{end}, s_1) \bigstep s_2) =
  ((\text{if} \ b \ s_1 \ \text{then} \ c_1 \ \text{else} \ c_2) \bigstep s_2)
  \]
\end{theorem}

In the previous theorems we have to be careful to not confuse the \enquote{if}
construct of the language \texttt{IMP} and that of the meta-language (in this case Lean).
To aid in this distinction we will use \texttt{monospace} font for the constructs
of the language, and normal font for the constructs of the meta-language.

\begin{theorem}\label{thm:while_rewrite}
  The \texttt{while} construct can be rewritten as:
  \begin{align*}
  &((\texttt{while} \ b \ \texttt{loop} \ c \ \texttt{end}, s_1) \bigstep s_3) = \\
  &\quad \text{if} \ b \ s_1 \\
  &\quad \text{then} \ \exists s_2, ((c, s_1) \bigstep s_2) \land ((\texttt{while} \ b \ \texttt{loop} \ c \ \texttt{end}, s_2) \bigstep s_3) \\
  &\quad \text{else} \ s_1 = s_3
  \end{align*}
\end{theorem}

\subsection{Behavioral equivalence}

The previous rewriting rules are useful in proofs as we can substitute one program
for another that is \enquote{exactly the same}, but we can also define a more
lax notion of equality that will let us say that two syntactically different
programs are \enquote{the same}. This is specially useful when proving
that certain optimizations don't change the behaviour of a program.

\begin{definition}[Behavioral equivalence]
  Two programs are behaviourally equivalent if they evaluate to the same state
  for all initial states:
  \[
  c_1 \sim c_2 \coloneqq \forall s_1 \ s_2, (c_1, s_1) \bigstep s_2 \iff (c_2, s_1) \bigstep s_2
  \]
\end{definition}

Is easy to show that this is an equivalence relation, that is, it is reflexive,
symmetric, and transitive. Armed with this we can, for example, show that adding
an arbitrary number of skips to a program doesn't change its behaviour.

\begin{theorem}
  Adding \texttt{skip} to the left of a program doesn't change its behaviour:
  \[
  \texttt{skip} \texttt{;} c \sim c
  \]
\end{theorem}

With this we can \enquote{optimize} the \texttt{skip} away of a program, but only
on the left, to completely remove them without consequence we need the following.

\begin{theorem}
  Adding \texttt{skip} to the right of a program doesn't change its behaviour:
  \[
  c \texttt{;} \texttt{skip} \sim c
  \]
\end{theorem}

A more interesting equivalence that mirrors Theorem~\ref{thm:while_rewrite} is
loop unfolding, a common optimization in compilers, used when memory is cheaper
than computation.

\begin{theorem}
  The \texttt{while} construct can be unfolded:
  \[
  \texttt{while} \ b \ \texttt{loop} \ c \ \texttt{end} \sim 
  \texttt{if} \ b \ \texttt{then} \ c \texttt{;} \ \texttt{while} \ b \ \texttt{loop} \ c \ \texttt{end} \ \texttt{else} \ \texttt{skip}
  \]
\end{theorem}

\begin{proof}
  Unfolding the equivalence definition the problem is reduced to proving that
  \(\forall s_1 \ s_3\):
  \[
    (\texttt{while} \ b \ \texttt{loop} \ c \ \texttt{end}, s_1) \bigstep s_3 \iff
    (\texttt{if} \ b \ \texttt{then} \ c \texttt{;} \ \texttt{while} \ b \ \texttt{loop} \ c \ \texttt{end} \ \texttt{else} \ \texttt{skip}, s_1) \bigstep s_3
  \]
  Now we can apply Theorem~\ref{thm:if_rewrite_2} to the right side of the equivalence:
  \[
    (\texttt{while} \ b \ \texttt{loop} \ c \ \texttt{end}, s_1) \bigstep s_3 \iff
    (\text{if} \ b \ s_1 \ \text{then} \ c; \texttt{while} \ b \ \texttt{loop} \ c \ \texttt{end} \ \text{else} \ \texttt{skip}, s_1) \bigstep s_3
  \]
  Now it is again a question of proving both implications separately.
  \begin{itemize}
    \item[(\(\implies\))] It suffices to check both possible constructors for
    the \texttt{while} construct in the left of the implication.
    \begin{itemize}
      \item[(\(b \ s_1\))]
      If the condition is true then the problem simplifies to:
      \[
        c; \texttt{while} \ b \ \texttt{loop} \ c \ \texttt{end}, s_1 \bigstep s_3
      \]
      and since we know the precondition is true we can apply the \(\texttt{whileTrue}\) constructor
      to obtain the derivation tree from the left side of the equivalence.
      \[
        \texttt{whileTrue}: \frac{
          b \ s_1 = \texttt{true}
          \quad
          (c, s_1) \bigstep s_2 
          \quad 
          (\texttt{while} \ b \ \texttt{loop} \ c \ \texttt{end}, s_2) \bigstep s_3
        }{
          (\texttt{while} \ b \ \texttt{loop} \ c \ \texttt{end}, s_1) \bigstep s_3
        }
      \]
      Now it suffices to apply the newfound hypotheses of \(\texttt{whileTrue}\) to 
      obtain the result:
      \[
        \texttt{cat}: \frac{
          (c, s_1) \bigstep s_2
          \quad
          (\texttt{while} \ b \ \texttt{loop} \ c \ \texttt{end}, s_2) \bigstep s_3
        }{
          (c; \texttt{while} \ b \ \texttt{loop} \ c \ \texttt{end}, s_1) \bigstep s_3
        }
      \]
      \item[(\(\lnot b \ s_1\))] The problem simplifies to:
      \[
        \texttt{skip}, s_1 \bigstep s_3
      \]
      proceeding in the same manner as before we can apply the \(\texttt{whileFalse}\) constructor:
      \[
        \texttt{whileFalse}: \frac{
          b \ s_1 = \texttt{false}          
        }{
          (\texttt{while} \ b \ \texttt{loop} \ c \ \texttt{end}, s_1) \bigstep s_1
        }
      \]
      so \(s_1 = s_3\) and we have the result.
    \end{itemize}

    \item[(\(\impliedby\))] This direction is also proved by cases on the condition 
    of the \texttt{if} construct, but is even easier than before, since we
    can directly simplify the precondition.
  \end{itemize}
\end{proof}

\subsection{Non termination}

When we attempted to define a \funfont{eval} function to evaluate commands before
we said that some programs could not end, what we mean by this is that the
big step relationship doesn't relate the program in some state to any state.

\begin{definition}[Non termination]
  We will say that a program doesn't terminate in a given \(s\) state if there 
  is no state that the program can evaluate to, that is:
  \[
   \forall s', (c, s) \centernot \bigstep s'
  \]
\end{definition}

We will prove that there are programs in \texttt{IMP} that don't terminate,
we will show that there is in fact a program that doesn't terminate in any
state.

\begin{theorem}\label{thm:non_termination}
  The program \texttt{while b loop c end}, where \(b \sim \texttt{true}\)
  doesn't terminate in any state.
  \[
  \forall s_1 \ s_2, (\texttt{while} \ b \ \texttt{loop} \ c \ \texttt{end}, s_1) \centernot \bigstep s_2
  \]
\end{theorem}

\begin{proof}
  In reality the \(\centernot \bigstep\) is syntactic sugar for
  the negation of the big step relation, so the theorem is equivalent to:
  \[
  (\texttt{while} \ b \ \texttt{loop} \ c \ \texttt{end}, s_1) \bigstep s_2 \implies \bot
  \]
  Now we have to use a little trick, we will first generalize the configuration on the
  left, calling it 
  \[\text{hconf} \coloneqq (\texttt{while} \ b \ \texttt{loop} \ c \ \texttt{end}, s_1)\] 
  with this we can rewrite the theorem as:
  \[
  \text{hconf} \bigstep s_2 \implies \bot
  \]
  now we will use the induction principle on the configuration, that is, we will
  prove the theorem for all the possible forms of \text{hconf}. Since
  \(\text{hconf} = (\texttt{skip}, s_1)\) for example is clearly false, we will only
  prove the case for the two constructors of \texttt{while}.

  In the \(\texttt{whileTrue}\) case we have to prove that:
  \[
    (\texttt{while} \ b \ \texttt{loop} \ c \ \texttt{end}, s_1) = (\texttt{while} \ b' \ \texttt{loop} \ c' \ \texttt{end}, s_1')
  \]
  but the induction hypothesis takes the form:
  \[
    \forall x, (\texttt{while} \ b \ \texttt{loop} \ c \ \texttt{end}, x) = (\texttt{while} \ b' \ \texttt{loop} \ c' \ \texttt{end}, s_2')
  \]
  so it suffices to choose \(x = s_2\). For the \(\texttt{whileFalse}\) case we 
  don't even need to use induction, in fact if \text{hconf} takes this form
  that means that the condition of the \texttt{while} construct is false but
  \(b \sim \texttt{true}\) by hypothesis so we have a contradiction.
\end{proof}

\subsection{Determinism}

Another important property of \texttt{IMP} that we can prove using the big step
semantics is that the evaluation of a program is deterministic, that is, given
a program and a state there is only one state that the program can evaluate to.

\begin{theorem}[\texttt{deterministic}]\label{thm:deterministic}
  The evaluation of a program is deterministic:
  \[
  \forall c \ s \ s_1 \ s_2, ((c, s) \bigstep s_1 \land (c, s) \bigstep s_2) \implies s_1 = s_2
  \]
\end{theorem}

\begin{proof}
  We will prove the theorem by induction on \((c, s) \bigstep s_1\).
  The \texttt{skip} and assignment cases are immediate, for the concatenation
  we will need the induction hypotheses:
  \begin{align*}
    \forall x, (c_1', s_1') &\bigstep x \implies s_3 = x \\
    \forall y, (c_2', s_2') &\bigstep y \implies s_1 = y
  \end{align*}
  the right hand side of the conjunction takes the form 
  \((c_1';c_2', s_1') \bigstep s_2\) so we must have its preconditions,
  that is:
  \[
    \texttt{ass}: \frac{
      (c_1', s_1') \bigstep s_2' 
      \quad 
      (c_2', s_2') \bigstep s_2
    }{
      (c_1';c_2', s_1') \bigstep s_2
    }
  \]
  Now we can apply the first induction hypothesis on the left precondition to
  obtain \(s_3 = s_2'\), with this the second precondition becomes
  \((c_2', s_3) \bigstep s_2\) and applying the second 
  induction hypothesis \(s_1 = s_2\), which is what we wanted to prove.

  The rest of the cases are analogous.
\end{proof}

\section{Small-Step Semantics}

Another way of defining the semantics of the language are small-step semantics. 
In this case instead of yielding directly the value of the expression we take one \enquote{single step}
that relates two configurations, with the state before taking the step and the 
state after, by doing this we can observe
the \emph{intermediate states} of the execution of a program, so every single
step of this semantics is \enquote{atomic}. This is very useful when we want to
prove properties of parallel or concurrent programs, although we won't go into
that in this text.
\begin{definition}[Small-step]
  We will denote the small-step relation \(\text{SmallStep}: \tyfont{Com} \times \tyfont{State}
\to \tyfont{Com} \times \tyfont{State} \to \Prop\) as:
  \[
    (c_1, s_1) \rightsquigarrow (c_2, s_2) \coloneq \text{SmallStep} \ (c_1, s_1) \ (c_2, s_2)
  \]
\end{definition}

The constructors of the relation are then:
\[
\begin{gathered}
  \texttt{ass}: (v \ \texttt{:=} \ a, s) \rightsquigarrow (\texttt{skip}, s[v \leftarrow s]) \quad
  \texttt{skipCat}: (\texttt{skip}; c, s) \rightsquigarrow (c, s) \quad
  \texttt{cat}: \frac{(c, s) \rightsquigarrow (c', s')}{(c; c'', s) \rightsquigarrow (c', c'', s')} \\
  \texttt{ifElse}: (\texttt{if} \ b \ \texttt{then} \ c \ \texttt{else} \ c' \ \texttt{end}, s)
    \rightsquigarrow (\text{if} \ b \ s \ \text{then} \ c \ \text{else} \ c', s) \\
  \texttt{whileLoop}: (\texttt{while} \ b \ \texttt{loop} \ c \ \texttt{end}, s)
    \rightsquigarrow (\text{if} \ b \ s \ \text{then} \ c;\texttt{while} \ b \ \texttt{loop} \ c \ \texttt{end}  \ \text{else} \ \texttt{skip}, s) 
\end{gathered}
\]
These rules capture the evaluation of commands from left to right in a sequential
fashion. Of course, we could have defined the \enquote{if} and \enquote{while}
constructs in other ways depending on the \enquote{granularity} that we want
in the semantics, evaluating the boolean could be a state and selecting the branch
another, for example. 

Let's see an example of how the small-step semantics work, for this we evaluate a
simple program:
\begin{example}\label{ex:small_step}
  Consider the program \texttt{x := 1; while x <= 1 loop x := x + 1 end},
  to ease reading we will denote the \enquote{while} block as \(W\).
  We have the evaluation tree:
  \[
  \texttt{cat}: \frac{
    \texttt{ass}: (\texttt{x := 1} , s_0) \rightsquigarrow (\texttt{skip}, s_0[x \leftarrow 1])
  }
  {(\texttt{x := 1;} W, s_0)
  \rightsquigarrow
  (\texttt{skip;} W, s_0[x \leftarrow 1])
  }
  \]
\end{example}

In this case the computation ends when we can't apply any rule to the program, 
that is when we have a configuration of the form \((\texttt{skip}, \cdot)\).
We will define the execution of a program from beginning to end with the small-step
semantics by using its \emph{reflexive transitive closure}.

\subsection{Reflexive transitive closure}

The reflexive transitive closure takes a binary relation and gives us he 
\enquote{minimal extension} of this relation that is \emph{reflexive}
and \emph{transitive}. Formally we can define it as a function that 
takes a binary relationship and return another.

\begin{definition}[Reflexive transitive closure]
  The reflexive transitive closure of a binary relation \(\mathcal{R}\) is the smallest 
  relation that contains \(\mathcal{R}\) and is both reflexive and transitive. 
  We can construct a relation \(\mathcal{R}^*\) from \(\mathcal{R}\) as an inductive type
  \(\tyfont{ReflTrans}: (\alpha \to \alpha \to \Prop) \to
  (\alpha \to \alpha \to \Prop)\) and denote it as 
  \(\mathcal{R}^* \coloneqq \tyfont{ReflTrans} \ \mathcal{R}\).

  The type \(\tyfont{ReflTrans}\) has two constructors:
  \[
    \texttt{refl}: a \mathcal{R}^* a
    \qquad 
    \texttt{tail}: \frac{
      a \mathcal{R}^* b
      \quad b \mathcal{R} c
    }{a \mathcal{R}^* c}
  \]
\end{definition}

As we can see this relation is not transitive by definition, but it is easy
to prove by induction.

\begin{theorem}
  The reflexive transitive closure is transitive, that is given 
  \(x \mathcal{R}^* y\) and \(y \mathcal{R}^* z\) we have \(x \mathcal{R}^* z\).
\end{theorem}

Another immediate but important property is that the relation itself contained
in the reflexive transitive closure.

\begin{theorem}
  The relation is contained in its reflexive transitive closure, that is
  \(a \mathcal{R} b \implies x \mathcal{R}^* y\).
\end{theorem}

Of course, we can not only append to the \enquote{tail} of the relation, but also
push to the \enquote{head}.

\begin{theorem}
  We can push a relation to the head of its reflexive transitive closure, that is,
  if we have \(a \mathcal{R} b\) and \(b mathcal{R}^* c\) then
  \(a \mathcal{R}^* c\).
\end{theorem}

This will become important soon, as we can define a principle of induction
on the \emph{head} of the relation, note that the structural induction
on the constructors works only on the \emph{tail} of the relation.

\begin{theorem}
  We can perform induction on the \emph{head} of the relationship, with the base
  case being the reflexivity, and the inductive step being the transitivity on
  the \emph{head} of the relation.
\end{theorem}

Another important theorem that we will use is the \emph{lifting} of a reflexive
transitive closure via a function, that is, if we have a relation between two
elements of a set we can lift this relation to the image of a function.

\begin{theorem}\label{thm:rel_lift}
  We can lift a reflexive transitive \(x \mathcal{R}^* y\) closure through a function
  \(f: \alpha \to \beta\). 
  Suppose that one relationship implies the other: 
  \[h: \forall a \ b, a \mathcal{R} b \to (f \ a) \mathcal{P} (f \ b),\]
  then we have \((f \ a) \mathcal{P}^* (f \ b)\). We can represent this as the 
  diagram~\ref{fig:rel_lift}, where
  the hypothesis \(h\) is expressed by the fact that the diagram commutes.
\end{theorem}

\begin{figure}[htb]\label{fig:rel_lift}
  \centering
  \begin{tikzcd}
    \alpha \times \alpha \arrow[dd, "f \times f"] \arrow[dr, "\mathcal{R}^*"] & \\
     & \Prop \\
    \beta \times \beta \arrow[ur, "\mathcal{P}^*"] &
  \end{tikzcd}
  \caption{Lifting of a reflexive transitive closure through a function.
           In reality, instead of \(\alpha \times \alpha\)
           and \(\beta \times \beta\) we should use
           the \emph{exponential objects} \(\alpha^\alpha\) and 
            \(\beta^\beta\). 
  }
\end{figure}

\begin{proof}
  We proceed by induction on \(a \mathcal{R}^* b\), the base case 
  \(\texttt{refl}: a \mathcal{R}^* a\) implies that \(a = b\), so
  \((f \ a) \mathcal{P}^* (f \ a)\)
  which is immediate by the reflexivity of the relation. For the inductive step:
  \[
    \texttt{tail}: \frac{
      a \mathcal{R}^* b 
      \quad 
      b \mathcal{R} c
    }{a \mathcal{R}^* c}
  \]
  we have to prove that \((f \ a) \mathcal{P}^* (f \ c)\), applying the
  constructor:
  \[
    \funfont{tail}: \frac{
      \text{ih}: (f \ a) \mathcal{P}^* (f \ b)
      \quad
      h: \displaystyle{\frac{
        b \mathcal{R} c
      }{
        (f \ b) \mathcal{P} (f \ c)
      }}
    }{
      (f \ a) \mathcal{P}^* (f \ c)
    }
  \]
  where \(\text{ih}\) is the induction hypothesis, we get the result.
\end{proof}

Now we can define the reflexive transitive closure of the small-step semantics.

\begin{notation}
  We will denote the reflexive transitive closure of the small-step semantics as
  \(\rightsquigarrow^*\).
\end{notation}

Now we can revisit the example~\ref{ex:small_step} with the reflexive transitive
closure of the small-step semantics.

\begin{example}
  As before, we consider the program \texttt{x := 1; while x <= 1 loop x := x + 1 end},
  to ease reading we will denote the \enquote{while} block as \(W\). We have the 
  partial derivation tree:
  \[
    \funfont{head}: \frac{
      \funfont{cat}: \displaystyle{\frac{
        \funfont{ass}: (\texttt{x := 1} , s_0) \rightsquigarrow (\texttt{skip}, s_0[x \leftarrow 1])
      }
      {(\texttt{x := 1;} W, s_0)
      \rightsquigarrow
      (\texttt{skip;} W, s_0[x \leftarrow 1])
      }}
      \quad
      \funfont{head}: \displaystyle{\frac{
        \funfont{skipcat}: \texttt{skip}
        \quad
        \vdots
      }{
        (W, s_0[x \leftarrow 1]) \rightsquigarrow^* (W, s_0[x \leftarrow 1][x \leftarrow 2])
      }}
    }{
      (\texttt{x := 1;} W, s_0) \rightsquigarrow^* (\texttt{skip;} W, s_0[x \leftarrow 1][x \leftarrow 2])
    }
  \]
  By successive applications of the constructors we get the result.
\end{example}

We preset now some theorems that will be very useful when proving the equivalence of
the big and small-step semantics, paying attention to the constructors of the small-step
relation we can see that the only \enquote{special} one is \texttt{cat}, since
it is the only one that requires an argument. For this reason all the following
theorems will give us analogous in \(\rightsquigarrow^*\) to the 
big-step semantics.

\begin{theorem}
  If an expression evaluates to a state concatenating that expression with another
  has that state as an intermediate. That is if \((c, s) \rightsquigarrow^* (\texttt{skip}, s')\)
  then \((c;c', s) \rightsquigarrow^* (\texttt{skip};c', s')\).
\end{theorem}

\begin{proof}
  The proof follows immediately by theorem~\ref{thm:rel_lift}.
  The relations \(\mathcal{R}\) and \(\mathcal{P}\) are both the small-step semantics
  \(\rightsquigarrow\), and the function \(f\) is the concatenation of commands:
  \(f: \texttt{Com} \times \texttt{State} \to \texttt{Com} \times \texttt{State}\)
  with \(f \coloneqq \lambda (c, s) \mapsto (c; c', s)\). Our hypothesis \(h\)
  is the concatenation of the small-step semantics, that is:
  \[h: \Forall_{(c, s), (d, t)} (c, s) \rightsquigarrow (d, t) \to (c; c', s) \rightsquigarrow (d; c', t)
     \coloneqq \lambda x \mapsto \funfont{cat} \ x\]
  and finally \((c, s) \rightsquigarrow^* (\texttt{skip}, s')\) by hypothesis.
\end{proof}

\begin{theorem}
  Concatenating an expression that evaluates to an intermediate configuration to
  another that evaluates in that configuration gives us the configuration of the
  second expression. That is if \((c, s) \rightsquigarrow^* (\texttt{skip}, s')\)
  and \((c', s') \rightsquigarrow^* (\texttt{skip}, s'')\) then
  \((c;c', s) \rightsquigarrow^* (\texttt{skip}, s'')\).
\end{theorem}

\begin{theorem}
  Concatenating an evaluation that doesn't change the state with another
  evaluation is equivalent to evaluating the second expression. That is
  if \((c, s) \rightsquigarrow^* (\texttt{skip}, s)\) then
  \((c;c', s) \rightsquigarrow^* (c', s)\).
\end{theorem}

\section{Denotational semantics}

We follow the construction presented  in~\cite[Chapter~11]{semantics_with_isabelle} 
and define a \emph{relational} denotational semantics for the language \texttt{IMP}.

As we saw before, defining a function that evaluates a program is not trivial,
in fact this function can never be \emph{total} as there are programs that don't
terminate (Theorem~\ref{thm:non_termination}). The function then
has to necessarily be \emph{partial}, that is, it doesn't return a value for
all inputs. This is a problem as in type theory all our functions are
total. To circumvent this problem we could use the \texttt{Option} monad,
and return \texttt{none} when the program doesn't terminate:
\[
  \funfont{eval} : \tyfont{Com} \to \tyfont{State} \to \tyfont{Option} \ \tyfont{State}.
\]

Instead of this, we will use the classic definition of relations as sets
\(\tyfont{Set} \ (\alpha \times \beta)\), this will allow us to use the 
partial ordering of subsets to define the semantics of the language via
the \emph{Knaster-Tarski} fix-point theorem.

\subsection{The Knaster-Tarski fix-point theorem}

The \emph{Knaster-Tarski} fix-point theorem is a generalization of the
\emph{Banach} fix-point theorem, it states that every monotone function
on a complete lattice has a least fix-point. This will be essential when defining
the denotational semantics if the \texttt{while} construct.
We begin by defining a \emph{partial order}.

\begin{definition}
  A partial order is a binary relation (usually denoted by \(\leq\)) that 
  satisfies the following properties:
  \begin{itemize}
    \item Reflexivity: \(a \leq a\).
    \item Transitivity: If \(a \leq b\) and \(b \leq c\) then \(a \leq c\).
    \item Less-than iff less-than-or-equal: \(a < b \iff a \leq b \land b \not\leq a\). 
    \item Antisymmetry: If \(a \leq b\) and \(b \leq a\) then \(a = b\).
  \end{itemize}
\end{definition}

Now we consider a set of elements with a partial order, if every set has an
an \emph{infimum} we call it a \emph{complete lattice}. Usually a complete lattice 
is defined as a  partially ordered type in which every set has a \emph{least upper bound}
and a \emph{greatest lower bound}, but these definitions are interchangeable.

\begin{definition}
  A complete lattice is a partially ordered type that has an infimum in every set. 
  That is, exists a function \(\funfont{Inf}: \tyfont{Set} \ \alpha \to \alpha\)
  that returns an element of the set satisfying the following properties:
  \begin{itemize}
    \item \(\forall S, \forall x \in S, \funfont{Inf} \ s \leq x\).
    \item \(\forall S, (\forall y \in S, x \leq y) \implies x \leq \funfont{Inf} \ s\).
  \end{itemize}
\end{definition}

Additionally, this infimum is \emph{unique}.

\begin{theorem}
  The infimum of a set in a complete lattice is unique.
  That is if \(\funfont{Inf} \ s = a\) and \(\funfont{Inf} \ s = b\) then \(a = b\).
\end{theorem}

Now we need to recall the definition of subsets in~\ref{sec:sets},
as sets, with the inclusion relation, form a partial order.

\begin{theorem}
  Sets with the inclusion relation form a partial order. Where
  \(A \leq B \coloneqq A \subseteq B\), \(A < B \coloneqq A \subset B\).
  and the infimum of a set is the intersection of all the sets in the set,
  that is \(\funfont{Inf} \ S = \{x \mid \forall A \in S, x \in A\}\).
\end{theorem}

One could read the infimum of the previous theorem as:
\(\funfont{Inf} \ S = \bigcap_{A \in S} A\), but we will refrain from doing so
since the definition of the infimum is more general than the intersection of sets,
and the indexed intersection is a bit subtle to define.

% Appendices 

% appendix
\appendix

\section{Normalization in Lean}
\subsection{Omega combinator}

\section{Equality: rec, ndrec and subst}

\section{\texttt{IMP} Turing completeness}

\section{Future work}

Things to add:

\begin{itemize}
  \item Add formalization of \(\lambda\)-calculus and proofs.
  \item \(\Omega\)-chain partial orders.
  \item Compiler.
  \item Typed \texttt{IMP}.
  \item Hoare Logic.
  \item Abstract interpretation.
\end{itemize}

\section{Code}

% Bibliography

\bibliographystyle{alpha}
\bibliography{refs}


\end{document}
